\subsection{Random Variables}

\hrulefill

\vspace{5mm}

In many random experiments, we are not interested in the specific outcome of an experiment, but in some characteristic of
that outcome. If you play the lottery , for example, you would like the numbers chosen in the weekly drawing to match
as many of the number you marked on your ticket as possible. You do not particularly care which of your numbers are
matched, but rather how many matches you have.

\paragraph*{Definition}
For a given sample space of some experiment, a \textbf{random variable} is any rule that associates a number with 
each outcome in S. In mathematical language, a random variable is a function whose domain is the sample space S and whose
range is the set of real numbers. We usually denote random variables with capital letters such as X, Y, and Z.

\begin{align*}
    X: S &\to \mathbb{R} \\
    \exists s \in S \ni X(s) &= x
\end{align*}

\paragraph*{Random Variable Examples}
\begin{itemize}
    \item Sum of two dice
    \item Number of defects on a randomly selected piece of fabric.
    \item SAT score of a randomly selected student.
    \item Number of heads in 10 flips of a coin.
    \item The amount a company has to pay to an employee in pension.
\end{itemize}

\paragraph*{Types of Random Variables}
\begin{itemize}
    \item A \textbf{discrete random variable} is a random variable that can take on a finite or countably infinite
    number of possible values.
    \item A \textbf{continuous random variable} is a random variable that can take on infinitely many values corresponding
    to an interval on the real line.
\end{itemize}